\clearpage

\lehead[]{\normalfont\sffamily\hspace*{-2.00cm}\textcolor{white}{\colorbox{lightblue}{\makebox[1.60cm][r]{\thechapter}}}\hspace{0.17cm}\textcolor{lightblue}{\chaptertitle}}
\rohead[]{\textcolor{lightblue}{\chaptertitle}\normalfont\sffamily\hspace*{0.17cm}\textcolor{white}{\colorbox{lightblue}{\makebox[1.60cm][l]{\thechapter}}}\hspace{-2.00cm}}
%\chead[]{}
\rehead[]{\textcolor{lightblue}{AvHG, Inf, My}}
\lohead[]{\textcolor{lightblue}{AvHG, Inf, My}}

\lstset{style=myJava}

\section{Thread-Synchronisation}

\subsection{Beispiel}

Zwei Threads benutzen einen gemeinsamen Zähler. Dadurch, dass das Betriebssystem
die Ausführung des Codes an jeder beliebigen Stelle unterbrechen kann
(\emph{Race Condition}), entstehen jedoch Fehler: Die ausgegebenen Zahlen
entsprechen nicht der gewünschten Reihenfolge, und mitunter gehen sogar
Zählschritte verloren.

\subsubsection{Hauptklasse}

\begin{lstlisting}
public class Synchronisation {
    static int zaehler = 0;
    
    public static void main(final String[] args) {
        MeinThread thread1 = new MeinThread(1);
        MeinThread thread2 = new MeinThread(2);
        thread1.start();
        thread2.start();
    }
}
\end{lstlisting}

\subsubsection{Thread-Klasse}

\begin{lstlisting}
public class MeinThread extends Thread {
    private int threadNummer;

    public MeinThread(int threadNummer) {
        this.threadNummer = threadNummer;
    }

    @Override
    public void run() {
        for (int i = 0; i < 100; i++) {
            Synchronisation.zaehler++;
            int help = Synchronisation.zaehler;
            æ// Simulation einer aufwändigen Berechnung:
æ            for (long x = 0; x < 1000000; x++) {
                Math.sin(x);
            }
            System.out.println("Thread " + threadNummer + ": " + help);
        }
    }
}
\end{lstlisting}

\subsection{Lösung des Problems}

Anweisungen, während deren Ausführung eine Unterbrechung (\emph{Task-Wechsel}
durch den \emph{Scheduler} des Betriebssystems) zu unerwünschten Ergebnissen
führt, werden als \emph{kritische Anweisungen} (bzw.\ \emph{kritischer
Abschnitt}) bezeichnet.

Anweisungen bzw.\ Anweisungsfolgen sind dann kritisch, wenn auf eine gemeinsam
genutzte Ressource (etwa eine Variable, eine Datei, ein Gerät oder auch ein
Textfeld) nicht nur lesend, sondern auch schreibend zugegriffen wird.

Man kann einen Block mit kritischen Anweisungen vor gleichzeitigen Zugriffen
schützen, indem man die entsprechenden Anweisungen in einen
\emph{synchronisierten Block} einfasst. Dazu benutzt man das Schlüsselwort
\lstinline|synchronized|. Über ein \emph{Monitor-Objekt} (oft auch einfach
\emph{Monitor} genannt) wird der exklusive Zugriff (wechselseitiger Ausschluss)
verwaltet.

Als Monitor-Objekt ist grundsätzlich jedes Objekt geeignet, auf das alle
Threads zugreifen können. Ein Thread kann den synchronisierten Block nur
betreten, wenn kein anderer Thread den entsprechenden Monitor bereits belegt.

Syntax:

\begin{lstlisting}
synchronized (monitorObjekt) {
    æ// Anweisungen, die nicht gleichzeitig von mehreren Threads
    // ausgeführt werden dürfen
æ}
\end{lstlisting}

Im obigen Beispiel könnte man in der Klasse \myClass{Synchronisation} ein
gemeinsames Monitor-Objekt \lstinline|static final Object monitor = new Object();|
anlegen und den kritischen Abschnitt in \lstinline|run()| damit absichern:

\begin{lstlisting}
synchronized (Synchronisation.monitor) {
    Synchronisation.zaehler++;
    help = Synchronisation.zaehler;
}
\end{lstlisting}

Wichtig ist, dass alle Threads dasselbe Objekt als Monitor benutzen!

Achtung: Damit alle Threads zügig zum Zug kommen, sollten immer so wenig
Befehle wie möglich in einem \lstinline|synchronized|-Block stehen (die
aufwändige Berechnung gehört also \emph{nicht} mit in den
\lstinline|synchronized|-Block).
